\subsubsection{Segment Tree Lazy Add + Set + Sum}
\begin{lstlisting}[style=mystyle, language=C++]
// TC: O(log n) por operacion, O(n) en memoria
// usa: sumar o asignar valores en rangos y consultar sumas
// lazy_set = LLONG_MAX indica "no hay set pendiente" (sirve para set a 0)
struct SegmentTree {
	struct Node {
		int val = 0, lazy_add = 0;
		int lazy_set = LLONG_MAX;
		Node(int x=0) { val = x; lazy_add = 0; lazy_set = LLONG_MAX; }
	};

	int n;
	vector<Node> st;

	SegmentTree(const vector<int>& a) {
		n = a.size();
		while(__builtin_popcount(n) != 1) n++;
		st.resize(2*n);
		for(int i=0 ; i<(int)a.size() ; i++)
		st[n+i] = Node(a[i]);
		for(int i=n-1 ; i>=1; i--)
		st[i] = merge(st[2*i], st[2*i+1]);
	}

	Node merge(Node &left, Node &right) {
		Node ans;
		ans.val = left.val + right.val;
		return ans;
	}

	void push(int node, int l, int r) {
		if(st[node].lazy_set != LLONG_MAX) {
			st[node].val = st[node].lazy_set * (r-l+1);
			if(l != r) {
				st[node*2].lazy_set = st[node].lazy_set;
				st[node*2].lazy_add = 0;
				st[node*2+1].lazy_set = st[node].lazy_set;
				st[node*2+1].lazy_add = 0;
			}
			st[node].lazy_set = LLONG_MAX;
		}
		if(st[node].lazy_add != 0) {
			st[node].val += st[node].lazy_add * (r-l+1);
			if(l != r) {
				st[node*2].lazy_add += st[node].lazy_add;
				st[node*2+1].lazy_add += st[node].lazy_add;
			}
			st[node].lazy_add = 0;
		}
	}

	void update(int node, int l, int r, int ql, int qr, int val,
	int type) {
		push(node, l, r);
		if(ql <= l && r <= qr) {
			if(type == 1) st[node].lazy_add += val;
			else st[node].lazy_set = val;
			push(node, l, r);
			return;
		}
		if(r < ql || qr < l) return;
		int mid = l + (r-l)/2;
		update(node*2, l, mid, ql, qr, val, type);
		update(node*2+1, mid+1, r, ql, qr, val, type);
		st[node].val = st[node*2].val + st[node*2+1].val;
	}

	int query(int node, int l, int r, int ql, int qr) {
		push(node, l, r);
		if(ql <= l && r <= qr) return st[node].val;
		if(r < ql || qr < l) return 0;
		int mid = l + (r-l)/2;
		return query(node*2, l, mid, ql, qr)
		+ query(node*2+1, mid+1, r, ql, qr);
	}

	// Funciones publicas
	void add_range(int l, int r, int val) {
		update(1, 0, n-1, l, r, val, 1);
	}
	void set_range(int l, int r, int val) {
		update(1, 0, n-1, l, r, val, 2);
	}
	int sum(int l, int r) {
		return query(1, 0, n-1, l, r);
	}
};
int main(){
	SegmentTree st({1, 2, 3, 4, 5});
	st.add_range(0, 4, 2);        // [3,4,5,6,7]
	cout << st.sum(0, 4) << endl; // 25
	st.set_range(1, 3, 0);        // [3,0,0,0,7]
	st.add_range(0, 4, 1);        // [4,1,1,1,8]
	cout << st.sum(0, 4) << endl; // 15
	st.set_range(2, 4, 5);        // [4,1,5,5,5]
	cout << st.sum(2, 4) << endl; // 15
	return 0;
}
\end{lstlisting}
